De este trabajo se puede concluir que la implementación de los algoritmos influye directamente en su comportamiento, generando 
diferencias entre el rendimiento teórico y el observado experimentalmente. Esto fue especialmente notorio en el caso del algoritmo 
de Strassen, donde la creación de copias y el uso intensivo de recursión provocan un alto consumo de memoria y un tiempo de ejecución 
mayor al esperado. Asimismo, se observó que para valores pequeños de \(n\) todos los algoritmos presentan un desempeño aceptable. 
Sin embargo, a medida que \(n\)crece, las diferencias entre ellos se vuelven más evidentes, lo que resalta la importancia de considerar 
tanto la complejidad teórica como la implementación práctica al evaluar algoritmos.